\scanpage{98}

et l'isom $E\longrightarrow\deleted{\unsure{E}}A\otimes L$ respecte les décompositions en somme directe

\diagram{\begin{tikzcd}[ampersand replacement=\&,column sep=small]
\deleted{\unsure{A}}E={} \& \underline L\arrow[d,"{-\mathrm{id}_{L}}"] \& {}\oplus{} \& \underline{\mathcal O}.x_0\arrow[d,dash,"\mathrm{id}"]\\
A\otimes L\simeq{} \& \underline L \& {}\oplus{} \& \underline{\mathcal O}.
\end{tikzcd}}

Si donc $Q$ est donnée par

\[
Q(\underset{\substack{\rotatebox{90}{$\in$}\\\Gamma L}}{u}+\lambda x_0)
=u^2+\lambda\deleted{\unsure{x}}\langle u,\underset{\substack{\rotatebox{90}{$\in$}\\\Gamma L}}{b}\rangle
+\lambda^2\underset{\substack{\rotatebox{90}{$\in$}\\\Gamma L^{\otimes2}}}{c},
\]

le signe $-\mathrm{id}_{L}$ sur le diagramme précédent implique que l'on aura

\[
N(\underset{\substack{\rotatebox{90}{$\in$}\\\Gamma L^\vee}}{v}+\deleted{\unsure{=}}\lambda 1)
=\deleted{\unsure{v^2}}\lambda^2\deleted{\unsure{+}}-\lambda\langle v,b\rangle+\langle v^{2},c\rangle
\]

La formule

\[
\boxed{Q_N(\varphi_N(x))=-N(x)\delta_N}
\]

(les deux mbres sont des \unsure{3|S}-formes p.r. à l'argument $N$ variable\ldots) est bien sûr valable dans le contexte général d'un mod. loc. libre $A$ \inserted{de rang $2$} et d'une forme quadratique \deleted{\unsure{invol.}} $N$ dessus, sans \unsure{autre donnée|autres données}.

\figure{\begin{tikzpicture}\draw (0,0)--(.35,-.95);\end{tikzpicture}}

% archive sheet number: 97

